\documentclass{article}
\usepackage{graphicx} % Required for inserting images

\title{Take-Home Eksamen DM500 Efterår 2024}
\author{Sirinapha Ruamwong og Erik Mortensen}
\date{November 2024}

\begin{document}

\maketitle

\section{Reeksamen februar 2015 Opgave 1 }

I det følgende laver vi $U = {1,2,3...,15}$ være universet (unisersal set). 
\vspace{10pt}


Betragt de to mængder:


$A = {\{2n | n\in S\}}$  og  $   B = {\{3n + 2| n \in S\}}$ 
\vspace{11pt}

Hvor $S = \{1,2,3,4\}$
\vspace{11pt}

Angiv samtlige elementer i hver af følgende mængder
\vspace{11pt}

a) $A$
\vspace{10pt}

b) $B$
\vspace{10pt}

c) $A \cap B$
\vspace{10pt}

 d) $A \cup B$
\vspace{10pt}

 e) $A - B$
\vspace{10pt}

 f) $\overline{\rm A}$  
\vspace{11pt}

 \subsection{Svar:}
 \vspace{11pt}
 Først skal jeg finde ud af hvad elementerne er i sæt $A$ og $B$.
 \vspace{10pt}

 $A = {\{2n | n\in S\}}$ : For alle $n$ elementer i $A$ er lig med $2$ gange $n$ elementer 
 
 i $S$.
\vspace{10pt}

Hvis $n = 1$ i $S$ , så $n = 1  *  2 = 2$ i $A$
\vspace{10pt}

Hvis $n = 2$ i $S$ , så $n = 2  *  2 = 4$ i $A$ 
\vspace{10pt}

Hvis $n = 3$ i $S$ , så $n = 3  *  2 = 6$ i $A$
\vspace{10pt}

Hvis $n = 4$ i $S$ , så $n = 4  *  2 = 8$ i $A$
\vspace{10pt}

$   B = {\{3n + 2| n \in S\}}$ : For alle $n$ elementer i $B$ er lig med $3$ gange $n + 2$ i $S$.  
\vspace{10pt}

Hvis $n = 1$ i $S$ , så $n = 1  *  3 + 2 = 5$ i $B$
\vspace{10pt}

Hvis $n = 2$ i $S$ , så $n = 2  *  3 + 2 = 8$ i $B$
\vspace{10pt}

Hvis $n = 3$ i $S$ , så $n = 3  *  3 + 2 = 11$ i $B$
\vspace{12pt}

Hvis $n = 4$ i $S$ , så $n = 4 *  3 + 2 = 14$ i $B$
\vspace{10pt}
\vspace{10pt}

$S = \{1,2,3,4\}$
\vspace{10pt}

\vspace{11pt}

a) $A = \{2, 4, 6, 8\}$ 
\vspace{10pt}

b) $B = \{5, 8, 11, 14\}$ 
\vspace{10pt}

c) $A \cap B = $  De elementer som fremgår både i $A$ og $B$. I dette tilfælde fremgår 

kun element $\{8\}$ i både sæt $A$ og $B$. 

\vspace{10pt}
 d) $A \cup B = $ De elementer som ikke fremgår i begge sæt dvs. $$\{2, 4, 5, 6, 11, 14\}$$
\vspace{10pt}
 
 e) $A - B = $ Alle elementer som er i $A$ men ikke i $B$. $$A=\{2, 4, 6, 8\} - B=\{5, 8, 11, 14\} = \{2,4,6\}$$
 
\vspace{10pt}

 f) $\overline{\rm A}=$ Alle elementer i  univers  $U = {1,2,3...,15}$ som ikke er i $A$. 
 $$U=\{1, 2, 3, 4, 5, 6, 7, 8, 9, 10, 11, 12, 13, 14, 15\} - A =\{2, 4, 6, 8\}$$ $$\overline{\rm A}= \{1, 3, 5, 7, 9, 10, 11, 12, 13, 14, 15\}$$

 \vspace{12pt}
 \section{Reeksamen februar 2015 Opgave 2}
 \vspace{10pt}
 Løs følgende opgaver: 
 \vspace{10pt}
 
a) Hvilke af følgende udsagn er sande? 
 \vspace{10pt}

1. $ \forall x \in \mathbb{N }: \exists y \in \mathbb{N}: x < y $  \vspace{10pt}
 
2. $ \forall x \in \mathbb{N }: \exists !y \in \mathbb{N}: x < y $ \vspace{10pt}

3. $\exists y \in \mathbb{N}: \forall x \in \mathbb{N}: x < y$
\vspace{10pt}


b) Angiv en negering af udsagn 1. fra spørgsmål a). Negerings-operationen 

$(\neg)$ må ikke indgå i dit udsagn. 
\vspace{11pt}

\subsection{Svar:}
\vspace{10pt}

Løsning for a) og b). 
\vspace{10pt}

Opgave a) 
 \vspace{10pt}
 
1. $$ \forall x \in \mathbb{N }: \exists y \in \mathbb{N}: x < y $$
Betyder: For alle $x$ i $\mathbb{N}$, der eksisterer $y$ i $\mathbb{N}$ således at $x$ er større end $y$.
\vspace{10pt}

Denne erklæring er sandt, fordi i en set af naturlige tal, for hvilken som helst

$x$, kan man altid finde en $y$ der er større end $x$. 
\vspace{10pt}

2. $$\forall x \in \mathbb{N }: \exists !y \in \mathbb{N}: x < y$$
Betyder: For alle $x$ i $\mathbb{N}$, der eksisterer præcis én $y$ i $\mathbb{N}$, således at $x$ er større end $y$.
 \vspace{10pt} 

Denne er erklæring er ikke sandt, fordi der findes ikke præcis én naturligt 

nummer $y$ som er større end $x$. 
 \vspace{10pt}

3. 
$$\exists y \in \mathbb{N}: \forall x \in \mathbb{N}: x < y$$ 
Betyder: Der eksisterer en $y$ i $\mathbb{N}$, således at alle $x$ i $\mathbb{N}$ er større end $y$. 
\vspace{10pt}

Denne erklæring er falsk, fordi der er ikke nogle naturlige numre i $y$ der er 

større end alle andre naturlige numre. 
\vspace{12pt}
 
\vspace{12pt}
 
\vspace{12pt}
  
\vspace{12pt}
   
\vspace{12pt}
  
\vspace{12pt}

\vspace{12pt}

\vspace{12pt}

\vspace{12pt}

\vspace{12pt}

\vspace{12pt}

\vspace{12pt}

Opgave b) 
 \vspace{10pt}
$$ \forall x \in \mathbb{N }: \exists y \in \mathbb{N}: x < y $$ Betyder: For alle $x$ i $\mathbb{N}$, der existerer $y$ i $\mathbb{N}$ således at $x$ er større end $y$.
 \vspace{10pt}

 For at negere overenstående erklæring uden at benytte negering $(\neg)$, skal jeg 
 
 lave sætningen om, således at det betyder det modsatte. 
 \vspace{10pt}
 
For at opnå dette, vil jeg flippe for alle $(\forall)$, til der eksisterer $(\exists)$ og der 

eksisterer $(\exists)$ til for alle $(\forall)$. 
 \vspace{10pt}
 
Så negationen vil blive: 
 \vspace{10pt}
 $$\exists x \in \mathbb{N} : \forall y \in \mathbb{N} : x \geq y$$
Dette betyder: Der eksistere en $x$ i $\mathbb{N}$, således at for alle $y$ i $\mathbb{N}$, at $x$ er lig med eller større en $y$. 
  
\section{Eksamen Januar 2015}
De givne funktioner er:
\[
f(x) = x^2 + x + 1 \quad \text{og} \quad g(x) = 2x - 2.
\]

\subsection*{(a) Er \( f \) en bijektion?}
For at afgøre, om \( f \) er en bijektion, skal vi undersøge, om den er både injektiv og surjektiv.

\begin{itemize}
    \item \textbf{Injektivitet:} \( f(x) = x^2 + x + 1 \) er en parabel, der åbner opad (fordi koefficienten foran \( x^2 \) er positiv). En parabel er ikke injektiv på \( \mathbb{R} \), da den ikke er én-til-én (den fejler den horisontale linjetest).
    
    \item \textbf{Surjektivitet:} For at være surjektiv skal alle \( y \in \mathbb{R} \) kunne opnås som \( f(x) \) for et \( x \in \mathbb{R} \). Da \( f(x) = x^2 + x + 1 \) har en minimumsværdi og ikke dækker hele \( \mathbb{R} \), er \( f \) heller ikke surjektiv.
\end{itemize}

Da \( f \) hverken er injektiv eller surjektiv, er den ikke en bijektion.

\subsection*{(b) Har \( f \) en invers funktion?}
En funktion har en invers, hvis og kun hvis den er bijektiv. Da \( f \) ikke er bijektiv, har den ikke en invers funktion.

\subsection*{(c) Angiv \( f + g \)}
Funktionen \( f + g \) er givet ved:
\[
(f + g)(x) = f(x) + g(x).
\]
Vi indsætter udtrykkene for \( f(x) \) og \( g(x) \):
\[
f(x) + g(x) = (x^2 + x + 1) + (2x - 2) = x^2 + 3x - 1.
\]
Så:
\[
(f + g)(x) = x^2 + 3x - 1.
\]

\subsection*{(d) Angiv \( g \circ f \)}
Sammensætningen \( g \circ f \) betyder, at vi først anvender \( f \) og derefter \( g \):
\[
(g \circ f)(x) = g(f(x)).
\]
Vi indsætter \( f(x) = x^2 + x + 1 \) i \( g(x) \):
\[
g(f(x)) = g(x^2 + x + 1) = 2(x^2 + x + 1) - 2 = 2x^2 + 2x + 2 - 2 = 2x^2 + 2x.
\]
Så:
\[
(g \circ f)(x) = 2x^2 + 2x.
\]

\end{document}
